% =========================================================
% Properties of Relations
% =========================================================

\section{Properties of Relations}

% ---------------------------------------------------------
% TOOLKIT
% ---------------------------------------------------------

% ---------------------------------------------------------
% Individual property definitions
% ---------------------------------------------------------
\begin{definitionbox}{Definition (Reflexive Relation)}
\begin{definition}[Reflexive Relation]\label{def:reflexive}
A binary relation $R\subseteq A\times A$ is \emph{reflexive} on $A$ if every
element of the domain is related to itself:
\[
  \forall a\in A,\ (a,a)\in R.
\]
In infix notation, this is the condition $\forall a\in A,\ aRa$.
\end{definition}
\end{definitionbox}

\begin{remark*}[Standard quantified statement]
Reflexivity is a universal requirement over the whole domain:
\[
  \operatorname{Reflexive}_A(R)
  \iff
  \forall a\in A,\ (a,a)\in R.
\]
For first-order logic with equality, the corresponding logical baseline is
the identity axiom
\[
  \forall x\,(x=x),
\]
which forces equality to behave reflexively in every structure.
\end{remark*}

\begin{dependencies}
\begin{itemize}
  \item \hyperref[def:relation]{Relation}
  \item \hyperref[def:universal-q]{Universal Formula}
\end{itemize}
\end{dependencies}

\begin{definitionbox}{Definition (Irreflexive Relation)}
\begin{definition}[Irreflexive Relation]\label{def:irreflexive}
A binary relation $R\subseteq A\times A$ is \emph{irreflexive} on $A$ if no
element of the domain is related to itself:
\[
  \forall a\in A,\ (a,a)\notin R.
\]
In infix notation, this is the condition $\forall a\in A,\ \neg(aRa)$.
\end{definition}
\end{definitionbox}

\begin{remark*}[Standard quantified statement]
Irreflexivity is not the same as merely failing to be reflexive:
\[
  \operatorname{Irreflexive}_A(R)
  \iff
  \forall a\in A,\ \neg(aRa),
\]
whereas
\[
  \neg\operatorname{Reflexive}_A(R)
  \iff
  \exists a\in A,\ \neg(aRa).
\]
Thus a relation may be neither reflexive nor irreflexive. In first-order
axiomatizations of strict order, irreflexivity is commonly imposed by the
axiom
\[
  \forall x\,\neg(x<x).
\]
\end{remark*}

\begin{dependencies}
\begin{itemize}
  \item \hyperref[def:relation]{Relation}
  \item \hyperref[def:universal-q]{Universal Formula}
\end{itemize}
\end{dependencies}

\begin{remark*}[Independence of reflexive and irreflexive]
A relation may be neither reflexive nor irreflexive (if some but not all
elements satisfy $(a,a)\in R$). However, a relation cannot be both reflexive
and irreflexive unless $A = \varnothing$.
\end{remark*}

\begin{definitionbox}{Definition (Symmetric Relation)}
\begin{definition}[Symmetric Relation]\label{def:symmetric}
A binary relation $R\subseteq A\times A$ is \emph{symmetric} on $A$ if every
related pair also appears in the reverse order:
\[
  \forall a,b\in A,\ (a,b)\in R \rightarrow (b,a)\in R.
\]
In infix notation, this is $\forall a,b\in A,\ aRb\rightarrow bRa$.
\end{definition}
\end{definitionbox}

\begin{remark*}[Standard quantified statement]
Symmetry is a conditional property:
\[
  \operatorname{Symmetric}_A(R)
  \iff
  \forall a,b\in A,\ (aRb\rightarrow bRa).
\]
It does not require every pair to be related; it requires only that any
existing directed relation can be reversed. For first-order logic with
equality, symmetry is expressed by the identity axiom
\[
  \forall x\forall y\,(x=y\rightarrow y=x).
\]
\end{remark*}

\begin{dependencies}
\begin{itemize}
  \item \hyperref[def:relation]{Relation}
  \item \hyperref[def:universal-q]{Universal Formula}
\end{itemize}
\end{dependencies}

\begin{definitionbox}{Definition (Antisymmetric Relation)}
\begin{definition}[Antisymmetric Relation]\label{def:antisymmetric}
A binary relation $R\subseteq A\times A$ is \emph{antisymmetric} on $A$ if any
two elements related in both directions must be equal:
\[
  \forall a,b\in A,\;
  \bigl((a,b)\in R\land(b,a)\in R\bigr)\rightarrow a=b.
\]
In infix notation, this is
$\forall a,b\in A,\ ((aRb\land bRa)\rightarrow a=b)$.
\end{definition}
\end{definitionbox}

\begin{remark*}[Standard quantified statement]
Antisymmetry forbids two-way relation cycles between distinct elements:
\[
  \operatorname{Antisymmetric}_A(R)
  \iff
  \forall a,b\in A,\ ((aRb\land bRa)\rightarrow a=b).
\]
It is not the same as asymmetry. Antisymmetry permits $aRa$; asymmetry forbids
it. Thus $\leq$ is antisymmetric, while $<$ is asymmetric. In first-order
axiomatizations of partial orders, antisymmetry is usually written
\[
  \forall x\forall y\,((x\leq y\land y\leq x)\rightarrow x=y).
\]
\end{remark*}

\begin{dependencies}
\begin{itemize}
  \item \hyperref[def:relation]{Relation}
  \item \hyperref[def:universal-q]{Universal Formula}
\end{itemize}
\end{dependencies}

\begin{definitionbox}{Definition (Asymmetric Relation)}
\begin{definition}[Asymmetric Relation]\label{def:asymmetric}
A binary relation $R\subseteq A\times A$ is \emph{asymmetric} on $A$ if every
related pair excludes the reverse pair:
\[
  \forall a,b\in A,\;
  (a,b)\in R \rightarrow (b,a)\notin R.
\]
In infix notation, this is
$\forall a,b\in A,\ (aRb\rightarrow\neg(bRa))$.
\end{definition}
\end{definitionbox}

\begin{remark*}[Standard quantified statement]
Asymmetry is the strict one-way condition
\[
  \operatorname{Asymmetric}_A(R)
  \iff
  \forall a,b\in A,\ (aRb\rightarrow\neg(bRa)).
\]
It implies irreflexivity by taking $a=b$, and it implies antisymmetry because
two-way relation between distinct elements is impossible. Conversely, a
relation that is both irreflexive and antisymmetric is asymmetric. For strict
orders, the typical first-order axiom is
\[
  \forall x\forall y\,(x<y\rightarrow\neg(y<x)).
\]
\end{remark*}

\begin{dependencies}
\begin{itemize}
  \item \hyperref[def:relation]{Relation}
  \item \hyperref[def:universal-q]{Universal Formula}
\end{itemize}
\end{dependencies}

\begin{remark*}[Asymmetric vs.\ antisymmetric]
Every asymmetric relation is antisymmetric, but not conversely. $\leq$ on
$\mathbb{R}$ is antisymmetric but not asymmetric (since $1 \leq 1$). Strict
order $<$ is both asymmetric and irreflexive.
\end{remark*}

\begin{definitionbox}{Definition (Transitive Relation)}
\begin{definition}[Transitive Relation]\label{def:transitive}
A binary relation $R\subseteq A\times A$ is \emph{transitive} on $A$ if any
two composable related pairs collapse to a direct related pair:
\[
  \forall a,b,c\in A,\;
  \bigl((a,b)\in R\land(b,c)\in R\bigr)\rightarrow(a,c)\in R.
\]
In infix notation, this is
$\forall a,b,c\in A,\ ((aRb\land bRc)\rightarrow aRc)$.
\end{definition}
\end{definitionbox}

\begin{remark*}[Standard quantified statement]
Transitivity is the formal chaining law:
\[
  \operatorname{Transitive}_A(R)
  \iff
  \forall a,b,c\in A,\ ((aRb\land bRc)\rightarrow aRc).
\]
For equality, the corresponding first-order axiom is
\[
  \forall x\forall y\forall z\,((x=y\land y=z)\rightarrow x=z).
\]
Together with reflexivity and symmetry, transitivity gives equivalence
relations; together with reflexivity and antisymmetry, it gives partial
orders.
\end{remark*}

\begin{dependencies}
\begin{itemize}
  \item \hyperref[def:relation]{Relation}
  \item \hyperref[def:universal-q]{Universal Formula}
\end{itemize}
\end{dependencies}

\begin{definitionbox}{Definition (Total (Connex) Relation)}
\begin{definition}[Total (Connex) Relation]\label{def:total-rel}
$R$ is \emph{total} (or \emph{connex}) if every pair is comparable:
\[
\forall a,b \in A, \;
(a,b) \in R \lor (b,a) \in R.
\]
\end{definition}
\end{definitionbox}

\begin{dependencies}
\begin{itemize}
  \item \hyperref[def:relation]{Relation}
  \item \hyperref[def:universal-q]{Universal Formula}
\end{itemize}
\end{dependencies}

% ---------------------------------------------------------
% Structural classes
% ---------------------------------------------------------
\begin{definitionbox}{Definition (Equivalence Relation)}
\begin{definition}[Equivalence Relation]\label{def:equivalence-rel}
$R$ is an \emph{equivalence relation} on $A$ if it is reflexive, symmetric,
and transitive.
\end{definition}
\end{definitionbox}

\begin{remark*}[Standard quantified statement]
$R$ is an equivalence relation on $A$ iff:
\[
(\forall a \in A,\; (a,a)\in R)\;\wedge\;
(\forall a,b \in A,\; (a,b)\in R \to (b,a)\in R)\;\wedge\;
(\forall a,b,c \in A,\; (a,b)\in R \wedge (b,c)\in R \to (a,c)\in R).
\]
\end{remark*}

\begin{remark*}[Role in mathematics]
Equivalence relations axiomatize the idea of ``sameness up to a chosen
criterion.'' They partition sets into equivalence classes and are the basis
for quotient constructions throughout algebra, topology, and analysis.
\end{remark*}

\begin{definitionbox}{Definition (Preorder)}
\begin{definition}[Preorder]\label{def:preorder}
$R$ is a \emph{preorder} on $A$ if it is reflexive and transitive.
\end{definition}
\end{definitionbox}

\begin{remark*}[Standard quantified statement]
$R$ is a preorder on $A$ iff:
\[
(\forall a \in A,\; (a,a)\in R)\;\wedge\;
(\forall a,b,c \in A,\; (a,b)\in R \wedge (b,c)\in R \to (a,c)\in R).
\]
\end{remark*}

\begin{definitionbox}{Definition (Partial Order)}
\begin{definition}[Partial Order]\label{def:partial-order}
$R$ is a \emph{partial order} on $A$ if it is reflexive, antisymmetric, and
transitive.
\end{definition}
\end{definitionbox}

\begin{remark*}[Standard quantified statement]
$R$ is a partial order on $A$ iff:
\[
(\forall a \in A,\; (a,a)\in R)\;\wedge\;
(\forall a,b \in A,\; (a,b)\in R \wedge (b,a)\in R \to a=b)\;\wedge\;
(\forall a,b,c \in A,\; (a,b)\in R \wedge (b,c)\in R \to (a,c)\in R).
\]
\end{remark*}

\begin{dependencies}
\begin{itemize}
  \item \hyperref[def:reflexive]{Reflexive Relation}
  \item \hyperref[def:antisymmetric]{Antisymmetric Relation}
  \item \hyperref[def:transitive]{Transitive Relation}
\end{itemize}
\end{dependencies}

\begin{definitionbox}{Definition (Total Order)}
\begin{definition}[Total Order]\label{def:total-order}
$R$ is a \emph{total order} on $A$ if it is a partial order and is total.
\end{definition}
\end{definitionbox}

\begin{remark*}[Standard quantified statement]
$R$ is a total order on $A$ iff $R$ is a partial order on $A$ and additionally:
\[
\forall a, b \in A,\; (a,b)\in R \vee (b,a)\in R.
\]
\end{remark*}

\begin{dependencies}
\begin{itemize}
  \item \hyperref[def:partial-order]{Partial Order}
  \item \hyperref[def:total-rel]{Total Relation}
\end{itemize}
\end{dependencies}

\begin{remark*}[Using structural properties in proofs]
When $R$ is asserted to be an equivalence relation or partial order, this is
shorthand for the conjunction of its defining properties. In proofs, cite
only the specific component needed: ``since $R$ is a partial order,
antisymmetry implies \dots'' rather than restating all properties.
\end{remark*}

\begin{example*}[Using a structural property in a proof]
Let $R$ be an equivalence relation on $A$. For any $a,b \in A$,
$(a,b) \in R \Rightarrow [a] = [b]$, where $[a] = \{x \in A \mid (a,x) \in R\}$.
\end{example*}



\begin{remark*}[Properties are logically independent]
No basic property implies another in general. A relation may be transitive
without being reflexive, or symmetric without being transitive. The structural
classes are distinguished precisely by requiring specific combinations.
\end{remark*}
